\clearpage
\section{Demonstrarea pe un exemplu concret principiul de lucru}
În timpul proiectării jocului a fost stabilit că se va numi \textbf{Eden}
iar codul sursă este accesibil pe pagina de GitHub --
\textbf{\url{https://github.com/catalinplesu/Eden}}. 
La această etapă în repozitoria cu, codul sursă a fost adăugată o Licență de tipul
\textbf{MIT} care permite persoanelor să facă ce vor cu codul sură.
Instrumentele utilizate pentru crearea acestui joc au fost enumerate în
capitolul precedent însă în pot fi utilizate și alte instrumente cu care
sunteți mai familiari.

\subsection{Realizarea proiectului jocului}
La această etapă a fost scris documentul de proiectare al
jocului(din engleză Game Design Document) care poate fi inspectat în anexa
\ref{anex:gdd} sau pe pagina de GitHub. Acțiunea jocului va avea loc într-o
pădure generată procedural. Jucătorul va avea ca scop să colecteze cât mai
multe fructe într-o perioadă limitată.

\subsection{Crearea unui proiect nou}
Pentru a crea acest proiect am utilizat comenzile demonstrate în capitolul
\ref{section:cargonew} utilizând numele proiectului respectiv.
\begin{lstlisting}
cargo new Eden
\end{lstlisting}
Merită menționat că Cargo a inițiat o nouă repozitorie \textbf{git} și a creat
un fișier \textbf{.gitignore}, pe parcursul dezvoltării acestui proiect au fost
făcute mai multe commit-uri pentru posibilitatea de a reveni la codul mai vechi
în cazul apariției problemelor. în fișierul \textbf{Cargo.toml} din rădăcina au
fost adăugate dependențele care au fost specificate în capitolul precedent.

\subsection{Aplicarea cunoștințelor despre ECS}
În acest joc sunt prezente mai multe entități, în sistemul de mai jos sunt
generate un număr de entități care reprezintă niște portocale. Una dintre componentele acestei entități este Transform căreia îi este dată valori aleatorii pe axele X și Z
valoare pe axa Y fiind prestabilită.
\begin{lstlisting}
fn spawn_fruits( ...) { ...
    for _ in 0..world_settings.fruits {
        let x = rng.gen_range(-half_size, half_size);
        let y = rng.gen_range(-half_size, half_size);
        commands
            .spawn_bundle((
                Transform::from_xyz(x, 10., y).
                with_scale(Vec3::new(scale, scale, scale)),
                GlobalTransform::identity(),
            ))
            .insert(Tangerine)
            ... } }
\end{lstlisting}
Parametrul \emph{scale} poate modifica mărimea portocalelor, ele fiind mai mari
ca în viața reală, acest lucru poate fi văzut în figura
\ref{fig:gameplay-fruits}.

\subsection{Scrierea plugin-urilor}
Pentru gruparea logică a sistemelor și entităților au fost create mai multe
plugin-uri pentru a le menține cu ușurință, aceasta este lista
plugin-urilor existente la moment:
\begin{itemize}
    \item PlayerPlugin;
    \item WorldPlugin;
    \item FruitsPlugin;
    \item UserInterfacePlugin.
\end{itemize}
În PlayerPlugin sunt grupate toate sistemele care țin de controlerul
jucătorului și care îl amplasează în scenă. În WorldPlugin se generează terenul
și sunt amplasați copacii și tufișiurile. În FruitsPlugin sunt generate
fructele și se conține sistemul care detectează daca jucătorul a colectat vreun
fruct, listingul acestui plugin este mai jos. Ultimul plugin
UserInterfacePlugin afișează pe ecran scorul jucătorului și cât timp mai are la
dispoziție până la finisarea jocului.
\begin{lstlisting}
pub struct FruitsPlugin;
impl Plugin for FruitsPlugin {
    fn build(&self, app: &mut App) {
        app.add_startup_system(spawn_fruits);
        app.add_system(tangerine_detection);
    }
}
\end{lstlisting}

\subsection{Mutarea plugin-urilor în module separate}
Scrierea unui program considerabil într-un singur fișier este o practică rea,
acest joc nu este foarte mare dar totuși are peste 1000 ceea ce a influențat
separarea codului în mai multe fișiere. În Rust fișierele înafră de main.rs se
numesc module, iar la dorință daca un modul este prea mare poate fi creat o
mapă care reprezină numele modulului și obligatoriu trebuie să conțină un
fișier numit mod.rs. Am decis să creez module pentru jucător, lume și interfața
de utilizator și să mut plugin-urile respective în ele. Structura proiectului
poate fi observată în figura \ref{fig:project_tree}.
\begin{figure}[H]
    \begin{center}
        \includegraphics[height=7cm]{project_tree}
        \caption{Structura proiectului}
        \label{fig:project_tree}
    \end{center}
\end{figure}

\subsection{Generarea procedurală a terenului}
Pentru generarea terenului a fost folosită funcția Perlin Noise. Aceasta este una dintre multele funcții oferite de librăria noise care poate fi utilizată adăugând următoarea linie în secțiunea de dependențe din fișierul Cargo.toml:
\begin{lstlisting}
noise = "0.7.0"
\end{lstlisting}
Mai jos este amplasată o secvență din funcția create\_mesh care demonstrează cum poate fi folosită funcția \textbf{noise\_map} pentru a genera datele pe axa Y având ca date de intrare coordonatele pe axa X și Z.
\begin{lstlisting}
fn create_mesh(n: u32, min: f64, max: f64) -> (Mesh, ColliderShape) {
    let perlin = Perlin::new();
    let noise_map = PlaneMapBuilder::new(&perlin)
        .set_size(n as usize, n as usize)
        .set_x_bounds(min, max)
        .set_y_bounds(min, max)
        .build();
    ...
    for x in 0..n + 1 {
        for z in 0..n + 1 {
            let y = noise_map.get_value(x as usize, z as usize) as f32 * 10.;
            ...
            heights[(z as usize, x as usize)] = y;
        }
    }
    let collider_shape = ColliderShape::heightfield(heights, ground_size);
    ...
}
\end{lstlisting}
Rezultatul executării acestei funcții poate fi observat în figura \ref{fig:gameplay-terrain} harta nu este foarte mare din motivul posibilităților limitate ale laptopului pe care a fost dezvoltat jocul. 

Pentru posibilitatea de a modifica cu ușurință parametrii lumii a fost creat un Singleton care conține mai mulți parametri.
\begin{lstlisting}
impl Default for WorldSettings {
    fn default() -> Self {
        Self {
            size: 200.,
            plants: 100,
            plant_dominance_offset: 0.,
            plants_colider: true,
            fruits: 100,
        }
    }
}
\end{lstlisting}
După cum poate fi observat cu ajutorul acestui singleton putem modifica cu
ușurință dimensiunile terenului, numărul de plante, specifica dacă dorim mai
mulți copaci sau arbuști, dacă este nevoie ca plantele să aibă un colider,
dezactivarea căruia ar mări performanța jocului corespunzător cu numărul de
plante și numărul de fructe.

\subsection{Finisarea jocului}
La această etapă a apărut idea de a adăuga un timp limită în care jucătorul
poate culege fructe. Când cronometrul este activ jucătorul se poate deplasa ca
în figura \ref{fig:gameplay-veiw} iar când cronometrul expiră pe ecran este
afișat END și jucătorul nu se mai poate mișca, ca în figura
\ref{fig:gameplay-end}. În secvența următoare de cod este arătată crearea
cronometrului care este o resursă.
\begin{lstlisting}
pub struct Countdown {
    pub timer: Timer,
}
impl Countdown {
    pub fn new() -> Self {
        Self {
            timer: Timer::from_seconds(60.0, false),
} } }
\end{lstlisting}
Imaginile cu capturi de ecran din joc pot fi observate în anexa \ref{anex:gameplay}. O demonstrare cu modul de funcționare a jocului este disponibilă pe YouTube la următorul url \url{https://youtu.be/HZdUG32fqB8}. 
